\begin{enumerate}
	\item
	\begin{enumerate}
		\item Diminuer de 2\,\% revient à multiplier par $1 - \dfrac{2}{100} = 0,98$.
		
		$u_1 = 0,98 \times \np{60000} = \np{58800}$.

		On prévoit \np{58800} naissances en 2026.

		\item $v_1 = \np{40000} + 900 = \np{40900}$.

		On prévoit \np{40900} décès en 2026.
	\end{enumerate}

	\item
	\begin{enumerate}
		\item D'une année à la suivante, le nombre de naissances diminue de 2\,\% donc, pour tout entier naturel $n$ :
		\[u_{n+1} = u_n - \dfrac{2}{100}u_n = u_n(1 - 0,02) = 0,98u_n.\]

		\item On passe d'un terme au suivant en multipliant toujours par le même nombre $0,98$ : la suite
		$(u_n)$ est \textbf{géométrique} de raison $q = 0,98$ et de premier terme $u_0 = \np{60000}$.

		\item Pour tout entier naturel $n$, $u_n = u_0 \times q^n$.
		
		 Pour tout entier naturel $n$, $u_n = \np{60000} \times 0,98^n$.
	\end{enumerate}

	\item
	\begin{enumerate}
		\item Pour tout entier naturel $n$, $v_{n+1} = v_n + 900$. La suite $(v_n)$
		est \textbf{arithmétique} de raison $r = 900$ et de premier terme $v_0 = \np{40000}$.

		\item Pour tout entier naturel $n$,  $v_n = v_0 + nr$
		
		Pour tout entier naturel $n$, $v_n = \np{40000} + 900n$.
	\end{enumerate}

	\item
	\begin{enumerate}
		\item On saisit en C3 la formule \texttt{=C2+900}.

		\item La ligne 12 correspond à $n = 10$ :
		$v_{10} = \np{40000} + 900 \times 10 = \np{49000}$.

		La cellule C12 affiche \np{49000}.
	\end{enumerate}

	\item Le solde naturel de l'année $2025+n$ est $u_n - v_n$ ; il est négatif lorsque $u_n < v_n$.

		Dans la tableau donné, on remarque que pour $n = 10$, $u_{10} > v_{10}$ et pour $n=11$, $u_{11} < v_{11}$.
	
	Comme $(u_n)$ est décroissante et $(v_n)$ croissante, le
	solde reste ensuite négatif.

	Le solde naturel devient négatif à partir de l'année
	$2025 + 11 = 2036$.
\end{enumerate}
